% it/figures/data_pipeline/tab_ae_gap.tex — versione italiana
\begin{table}[!ht]
    \centering
    \small
    \caption[Autoencoder contro PCA a dimensione latente appaiata.]{Ricostruzione
    dei profili verticali da parte di un autoencoder non lineare contro la PCA lineare,
    a dimensione latente $k$ appaiata. Lo scarto è negativo a ogni $k$: la PCA non è mai
    battuta, quindi la varietà è piatta. Gli scarti sono calcolati da valori non
    arrotondati, perciò possono differire di 0.01 dalla differenza delle colonne
    arrotondate.}
    \label{tab:ae-gap}
    \begin{tabular}{@{}crrr@{}}
        \toprule
        \textbf{$k$} & \textbf{PCA cum.\ \%} & \textbf{AE $R^2_{\mathrm{Sv}}$ \%} & \textbf{Scarto (pp)} \\
        \midrule
        1 & 89.82 & 89.67 & $-0.15$ \\
        2 & 97.86 & 97.77 & $-0.08$ \\
        3 & 99.42 & 99.36 & $-0.06$ \\
        4 & 99.80 & 99.74 & $-0.06$ \\
        5 & 99.91 & 99.83 & $-0.08$ \\
        \bottomrule
    \end{tabular}
\end{table}
